% Source: https://artofproblemsolving.com/community/c3h2479782_conditional_statement
% Stable identity: aops-topic-2479782-q01
\textbf{(a)} A conditional statement is true only if the conclusion holds in \emph{every} case where the hypothesis holds. For $x=2$ the hypothesis $x^2=4$ is true but the conclusion $x=-2$ is false. So this counterexample shows that the statement is \textbf{false}.

\textbf{(b)} If $x=-2$, then $x^2=(-2)^2=4$, so the converse is \textbf{true}.

\textbf{(c)} From $x^2=4$ we get $(x-2)(x+2)=0$, so $x=2$ or $x=-2$. The true statement is: ``If $x^2=4$, then $x=2$ or $x=-2$.''

So a conditional whose conclusion lists only one of the two solutions is false.
